\clearpage
\SourceLeaf{328}
\pdfbookmark[1]{Bemerkungen zu den Zusätzen}{zusaetze-bemerkungen}
\begin{center}
\textbf{BEMERKUNGEN.}
\end{center}

Die Notizen [2], [4], [5], [13], [15] finden sich in einem Handbuche; [3], [9], [10], [17], [18] auf losen Blättern; [1] u. [12] in \textsc{Gauss}’ Exemplaren des Berliner Astronomischen Jahrbuchs für 1805 und 1818; [6] in \textsc{Gauss}’ Exemplar von \textsc{Vega}, log.-trig. Handbuch, Leipzig, 1797; [7], [8], [11], [16] sind bereits an den angegebenen Orten veröffentlicht. Die Notiz [14] ist einem Heft mit dem Titel »Abhandlungen von \textsc{Carl Friedrich Gauss}, Erstes Heft« entnommen, welches von Herrn Professor \textsc{Jacobi} in Göttingen freundlichst zur Verfügung gestellt wurde und aus dem Nachlass von \textsc{Wilhelm Weber} stammt; im Einbande findet sich die Namenseintragung »\textsc{J. H. Westphal}, Danzig 1818«; es enthält ausser einigen Copien \textsc{Gauss}scher Aufsätze Notizen derselben Art, wie die Notiz [14], welche aus Vorlesungen oder sonstigen mündlichen Mittheilungen von \textsc{Gauss} herzurühren scheinen.

Der zweite Theil der Notiz [6] ist mit aufgenommen worden, obwohl er in keinem direkten Zusammenhang mit den Art. 88---90 steht; man vergleiche Band VIII, S. 128.

Die Formeln der Notiz [15] finden sich bereits in der Disquisitio de elementis ellipticis Palladis (Bd. VI, S. 17---18); sie sind hier in eine für die numerische Rechnung geeignetere Form gebracht. Auch das Resultat des numerischen Beispiels findet man a. a. O. S. 19; die Abweichung beim Werthe von \(C\) rührt daher, dass in der Notiz [15] bei der Herausgabe ein Rechenfehler verbessert wurde.

Über die Grenzen der geocentrischen Örter der Himmelskörper (Zodiakus, Notiz [16]---[18]) finden sich im Nachlass auf losen Blättern noch mehr einzelne Entwickelungen und Rechnungen, deren Bearbeitung sich aber wegen ihrer allzu fragmentarischen Beschaffenheit als unthunlich erwies.

Unerhebliche Änderungen des Originaltextes haben an mehrern Stellen vorgenommen werden müssen.
\begin{flushright}
\textsc{Brendel}.
\end{flushright}
\begin{center}
\rule{0.3\linewidth}{0.4pt}
\end{center}
